\chapter{Pruebas}\label{cap:pruebas}

\section{Introducción}
En este capítulo explicaremos...

\section{Conclusiones}
En este capítulo concluimos que\dots

\section{Aseguramiento de Calidad y Entorno de Ejecución}
La inmutabilidad intrínseca de los contratos inteligentes en Ethereum hace que sea imposible aplicar parches de seguridad convencionales una vez el código ha sido desplegado. Por ello, el ecosistema TraDeck se ha desarrollado bajo una estricta metodología de pruebas (\textit{Testing}).

\subsection{Entorno de Pruebas Unitarias}
Para la simulación y validación de la lógica de negocio, se ha utilizado el \textit{framework} Hardhat en combinación con las librerías Ethers.js y Chai. Esto permitió levantar una red blockchain efímera en memoria y ejecutar pruebas de caja blanca sobre los contratos.

El conjunto de pruebas validó de forma matemática:
\begin{itemize}
    \item \textbf{Control de Acceso:} Verificación de que solo el propietario de un token ERC-721 puede invocar la función \texttt{listForSale}.
    \item \textbf{Aislamiento de Fondos:} Comprobación de que la función \texttt{buyCard} retiene exactamente la cantidad de Wei (la fracción más pequeña de ETH) especificada por el vendedor.
    \item \textbf{Reversión de Estado:} Uso de aserciones para confirmar que transacciones maliciosas (ej. intentar confirmar una entrega sin ser el comprador) disparan los mecanismos \texttt{require} y revierten el estado de la red sin consumir gas innecesario.
\end{itemize}